% ============================================================================
%  Appendix E: Derivations for Section 4
%  Proof of Lemma 1 (claiming-age accounting) and the algebra of the
%  decomposition in Section 4.5. No numbers.
% ============================================================================
\section{Derivations}\label{app:derivations}

This appendix proves Lemma~\ref{lem:claiming}, gives a sufficient condition under which the fair credit is higher in worse health, decomposes the change in retirement outlays caused by a change in the claiming-age schedule, and derives the relations between the net cost, the willingness to pay, and the fiscal multiplier used in Section~\ref{sec:fw_decomp}.

\subsection{Proof of Lemma~\ref{lem:claiming}}

Let $b(\cdot)$ be any schedule of benefits by claiming age and $\Delta b(a)\equiv b(a+1)-b(a)$ its increment at age $a$. The definition of $A(\cdot)$ implies $A(a)=(1+r)^{-(a-62)}\mathcal{S}(a)+A(a+1)$. Hence
\begin{equation*}
\begin{split}
B(a+1)-B(a)&=b(a+1)\,A(a+1)-b(a)\,A(a)\\
&=\Delta b(a)\,A(a+1)-b(a)\,(1+r)^{-(a-62)}\,\mathcal{S}(a).
\end{split}
\end{equation*}
For $a\geq\mathit{NRA}$, $\Delta b(a)=\delta\,\mathit{PIA}$, which gives \eqref{eq:claimnpv}. Since $a+1\leq 70\leq T$ and $\mathcal{S}(a+1)>0$, $\mathit{PIA}\,A(a+1)>0$, and $B(a+1)<B(a)$ if and only if $\delta<\delta^{*}(a)$ as defined in \eqref{eq:faircredit}. Dividing the numerator and the denominator of $\delta^{*}(a)$ by $(1+r)^{-(a-62)}\mathcal{S}(a)>0$ gives the second expression in \eqref{eq:faircredit}, $\delta^{*}(a)=[b(a)/\mathit{PIA}]/\ddot a(a)$, where
\begin{equation*}
\ddot a(a)\equiv\sum_{t=a+1}^{T}(1+r)^{-(t-a)}\,\frac{\mathcal{S}(t)}{\mathcal{S}(a)}
\end{equation*}
is the expected present value at age $a$, for an individual alive at $a$, of one dollar a year from age $a+1$ until death. The factor $b(a)/\mathit{PIA}=1+\delta(a-\mathit{NRA})$ does not depend on survival, and $\ddot a(a)$ is strictly increasing in each $\mathcal{S}(t)/\mathcal{S}(a)$, $t>a$, and nondecreasing in $T$, since a later terminal age adds nonnegative terms. Therefore $\delta^{*}(a)$ is strictly decreasing in conditional survival at every age $t>a$ and nonincreasing in $T$. \hfill$\square$

Three remarks qualify the lemma. First, the identity holds before the NRA with the corresponding increment: a one-year delay at any age $a$ lowers the present value of benefits if and only if $\Delta b(a)/\mathit{PIA}<[b(a)/\mathit{PIA}]/\ddot a(a)$, the benefit forgone divided by the annuity factor. Under the schedules of \autoref{tab:schedules}, $\Delta b(a)/\mathit{PIA}$ is $1/15$ in the three years before the NRA, $1/20$ in earlier years, $\delta$ after the NRA, and $\psi$ at every age under the symmetric reforms. Second, only conditional survival $\mathcal{S}(t)/\mathcal{S}(a)$ enters $\delta^{*}(a)$, so the lemma holds with survival probabilities conditional on any information available at age $a$, in particular on health at $a$. Third, the lemma abstracts from the earnings test and from changes in the PIA. Additional earnings after age 60 raise the AIME only if they exceed its current value \eqref{eq:aime}, and benefits withheld by the earnings test are credited as if claiming had been delayed (Section~\ref{sec:model_ss}); both add terms to the fiscal value of a delay by a worker who keeps working.

\subsection{Health and the Fair Credit}

Order health states from best to worst, $h\in\{0,1,2\}$, and write conditional survival from age $a$ to age $t>a$ as $\mathcal{S}(t)/\mathcal{S}(a)=\E\big[\prod_{k=a+1}^{t}S(h_{k-1},k)\,\big|\,h_a\big]$, where the expectation is over the health paths generated by the transition probabilities $\pi_{j,i,k}$. Suppose that (i) one-period survival $S(h,k)$ is nonincreasing in $h$ at every age, and (ii) health transitions are stochastically monotone: for every age $k$ and state $j$, $\Pr(h_{k}\geq j\mid h_{k-1}=i)$ is nondecreasing in $i$. Then an individual in health $h_a$ has weakly higher conditional survival at every horizon than an individual in health $h'_a>h_a$, and hence, by Lemma~\ref{lem:claiming}, a weakly lower fair credit $\delta^{*}(a)$.

To see this, construct both health paths from common uniform draws $U_k$, setting $h_k$ equal to the smallest $j$ such that $\Pr(h_k\leq j\mid h_{k-1})\geq U_k$. By (ii), the conditional distribution functions are ordered, so $h_{k-1}\leq h'_{k-1}$ implies $h_k\leq h'_k$; by induction, the path that starts in better health is never in worse health. By (i), each factor $S(h_{k-1},k)$ is then weakly larger along that path, and so is the product. Taking expectations gives the result. In the estimated model, one-period survival is lower in poor health and with a disability at every age (panel B of \autoref{fig:firststep}), which is condition (i). We do not verify condition (ii) for the estimated transition probabilities, of which panel A of \autoref{fig:firststep} shows only those from good health. % AUTHORS: check condition (ii) in the estimated transition matrices (first-step inputs only) and state the result.

\subsection{Retirement Outlays under a Change in the Claiming-Age Schedule}

Consider a population of $N$ individuals indexed by $i$, with primary insurance amounts $\mathit{PIA}_i$, survival probabilities $\mathcal{S}_i(\cdot)$, and annuity values $A_i(\cdot)$ as in Lemma~\ref{lem:claiming}. Under the status-quo schedule $b^0_i(\cdot)$, individual $i$ claims at age $a^0_i$; under the reformed schedule $b^1_i(\cdot)$, at age $a^1_i$. Write $B^{j}_i(a)=b^{j}_i(a)A_i(a)$. Holding each PIA fixed and abstracting from the earnings test, the change in the per capita present value of retirement outlays is, in expectation over survival,
\begin{equation*}
\Delta\mathit{SS}
=\underbrace{\frac{1}{N}\sum_{i}\big[b^{1}_i(a^{0}_i)-b^{0}_i(a^{0}_i)\big]A_i(a^{0}_i)}_{\text{mechanical}}
\;+\;\underbrace{\frac{1}{N}\sum_{i}\big[B^{1}_i(a^{1}_i)-B^{1}_i(a^{0}_i)\big]}_{\text{behavioral}}.
\end{equation*}
For an individual who delays, $a^1_i>a^0_i$, the behavioral term telescopes into one-year steps, $B^{1}_i(a^{1}_i)-B^{1}_i(a^{0}_i)=\sum_{a=a^0_i}^{a^1_i-1}\big[B^{1}_i(a+1)-B^{1}_i(a)\big]$, each valued by the identity of the proof of Lemma~\ref{lem:claiming} under the new schedule; for an individual who claims earlier, it is minus the corresponding sum from $a^1_i$ to $a^0_i-1$.

For an increase in the DRC from $\delta$ to $\delta+d\delta$, $b^{1}_i(a)-b^{0}_i(a)=(a-\mathit{NRA})\,\mathit{PIA}_i\,d\delta$ for $a>\mathit{NRA}$ and zero otherwise, so the mechanical term is
\begin{equation*}
\frac{1}{N}\sum_{i}\Ind\{a^{0}_i>\mathit{NRA}\}\,(a^{0}_i-\mathit{NRA})\,\mathit{PIA}_i\,A_i(a^{0}_i)\,d\delta\;\geq\;0,
\end{equation*}
the additional benefits of the inframarginal late claimers. Each one-year delay after the NRA contributes a negative amount to the behavioral term if and only if $\delta+d\delta<\delta^{*}_i(a)$, evaluated under the new schedule, and each delay before the NRA if and only if $\Delta b^{1}_i(a)<b^{1}_i(a)/\ddot a_i(a)$, where $\Delta b^{1}_i(a)$ is the increment of the new schedule at age $a$. The direct cost of the reform is negative if and only if the behavioral term is negative and larger in absolute value than the mechanical term. For an increase in the NRA, $b^{1}_i(a)<b^{0}_i(a)$ at every claiming age from 62 to 70 (\autoref{tab:schedules}), so the mechanical term is negative whatever the claiming ages, and behavioral responses add to or subtract from this saving. In both cases, the earnings test and the effect of additional earnings on the PIA add terms to $\Delta\mathit{SS}$, and the effects on taxes, DI, and transfers enter the indirect cost.

\subsection{Mechanical Effects, Fiscal Externalities, and Program Accounts}

Let $\zeta^{0}$ and $\zeta^{1}$ denote the rules of the status quo and of the reform, and $\mathcal{B}^{0}$ and $\mathcal{B}^{1}$ the simulated histories of the cohort under them: hours, claiming dates, and asset holdings at every age, together with the health, survival, and wage draws, which are the same under both. Let $\Gamma(\zeta,\mathcal{B})$ be the per capita present value at age 62 of the cost aggregate of the reform family---$G$, $G+\mathit{MC}^{tg}$, or $G+K$---computed by applying the rules $\zeta$ to the histories $\mathcal{B}$. Then
\begin{equation*}
\mathit{NC}=\Gamma(\zeta^{1},\mathcal{B}^{1})-\Gamma(\zeta^{0},\mathcal{B}^{0})
=\underbrace{\Gamma(\zeta^{1},\mathcal{B}^{0})-\Gamma(\zeta^{0},\mathcal{B}^{0})}_{\mathit{ME}}
+\underbrace{\Gamma(\zeta^{1},\mathcal{B}^{1})-\Gamma(\zeta^{1},\mathcal{B}^{0})}_{\mathit{FE}},
\end{equation*}
which is the first equality in \eqref{eq:netcost}: the mechanical effect applies the new rules to unchanged behavior, and the fiscal externality is the change in cost caused by the change in behavior under the new rules.

The cost aggregate is a sum of signed program accounts, $\Gamma=\sum_{p}\Gamma_{p}$, with outlays entering positively and revenues negatively. Let $o$ denote the program the reform changes. The direct cost is $\mathit{DC}=\Delta\Gamma_{o}=\mathit{ME}_{o}+\mathit{FE}_{o}$ and the indirect cost $\mathit{IC}=\sum_{p\neq o}\Delta\Gamma_{p}=\mathit{ME}_{-o}+\mathit{FE}_{-o}$, with each term defined as above for the corresponding accounts. Summing gives the second equality in \eqref{eq:netcost}. The mechanical effect on other programs, $\mathit{ME}_{-o}$, is nonzero when programs interact at given behavior. For example, the consumption floor \eqref{eq:transfer} falls one for one with benefits for an individual whose transfer is positive both before and after the reform, so a benefit increase lowers transfers mechanically.

For a small reform, the fiscal externality can be attributed to margins to first order: it is the sum, over the choices that the reform changes, of the induced change in each choice times its fiscal value, the change in the cost aggregate $\Gamma$ that the choice causes when the individual's other choices are held fixed. The fiscal value of a one-year delay of the claim is the change in outlays of Lemma~\ref{lem:claiming}, together with the changes in transfers and in benefits withheld by the earnings test. The fiscal value of participation at age $t$ is minus the participation tax,
\begin{equation*}
\tau^{P}_t=\big[T_t(y^{1}_t)-T_t(y^{0}_t)\big]+T^{p}_t(w_t n_t)+\big(db^{0}_t-db^{1}_t\big)+\big(tr^{0}_t-tr^{1}_t\big)+\big(ss^{0}_t-ss^{1}_t\big)-\Delta B^{aime}_t,
\end{equation*}
where superscripts 1 and 0 denote values when the individual works $n_t$ hours and when she does not work, the term in $ss_t$ is the withholding by the earnings test of a beneficiary, and $\Delta B^{aime}_t$ is the present value of the increase in future retirement and disability benefits caused by the year's earnings, through the AIME and through the crediting of benefits withheld by the earnings test \citep[see][for participation tax rates]{kleven2006marginal}. The fiscal value of an additional dollar of assets is minus the present value of the income taxes on the additional interest income and of the reduction in means-tested transfers. For nonmarginal reforms, such an attribution depends on the order in which margins are changed; the program accounts do not.

\subsection{The Willingness to Pay for a Small Reform}

Index the rules by a scalar $\zeta$ that affects only the resources available at each age,
\begin{equation*}
R_t\equiv a_t+Y_t+b_t\,ss_t+db_t+tr_t-m_t,
\end{equation*}
so that $c_t=R_t-a_{t+1}$ by \eqref{eq:budget}; every reform of Section~\ref{sec:framework} has this form, since benefits, taxes, out-of-pocket medical expenses, and the return on assets enter only $R_t$. Treat hours $n_t$, claiming $b_t$, and next-period assets $a_{t+1}$ as the choices. Holding them fixed, a change in $\zeta$ changes consumption one for one with resources, $dc_t=(\partial R_t/\partial\zeta)\,d\zeta$, where $\partial R_t/\partial\zeta$ is the mechanical change in resources at age $t$, including the offsetting change in consumption-floor transfers. Because individuals choose optimally, the envelope theorem gives
\begin{equation*}
\frac{dV_{30}(X_{30})}{d\zeta}=\E_{30}\sum_{t=30}^{T}\beta^{t-30}\Bigg(\prod_{k=31}^{t}s_k\Bigg)\,u_c(c_t,l_t)\,\frac{\partial R_t}{\partial\zeta},
\end{equation*}
evaluated at the optimal choices. The theorem applies with discrete participation and claiming choices because $V_{30}$ is the upper envelope of the values of the discrete alternatives and is differentiable at almost every $\zeta$ \citep{milgrom2002envelope}; individuals who switch alternatives in response to the reform are indifferent between them at the margin, so their switch has no first-order effect on welfare.

Homogeneity of \eqref{eq:utility} implies $u_c(c,l)\,c=\gamma(1-\nu)\,u(c,l)$, so that $\gamma(1-\nu)\,\bar V^{0}$ equals the average over individuals of $\E_{30}\sum_{t}\beta^{t-30}(\prod_{k}s_k)\,u_c(c_t,l_t)\,c_t$. Differentiating \eqref{eq:ce} at the status quo, the CE of a small reform $d\zeta$ is
\begin{equation*}
d\,\mathit{CE}=\frac{d\bar V/d\zeta}{\gamma(1-\nu)\,\bar V^{0}}\,d\zeta
=\frac{\overline{\E_{30}\sum_{t}\beta^{t-30}\big(\prod_{k}s_k\big)\,u_c(c_t,l_t)\,\partial R_t/\partial\zeta}}{\overline{\E_{30}\sum_{t}\beta^{t-30}\big(\prod_{k}s_k\big)\,u_c(c_t,l_t)\,c_t}}\;d\zeta,
\end{equation*}
where the bars denote averages over individuals. The numerator is the marginal-utility-weighted mechanical transfer, the willingness to pay of \citet{hendren2016policy} expressed in utility units; the denominator is the marginal-utility-weighted consumption stream, which converts it into a uniform proportional change in consumption. Behavioral responses do not enter. The dollar willingness to pay of \eqref{eq:mvpf}, $\WTP^{\$}=\bar C^{0}_{s}\,d\,\mathit{CE}$, and the value of consumption smoothing per mechanical dollar, $\WTP^{\$}/\mathit{ME}$, depend both on how the transfer is distributed across states and ages relative to consumption and on the base $\bar C^{0}_{s}$.

\subsection{The Fiscal Multiplier and the Desirability of a Reform}

Combine a reform with net cost $\mathit{NC}$ with a change in another policy that balances the budget: if $\mathit{NC}>0$, the other policy is cut so as to raise $\mathit{NC}$; if $\mathit{NC}<0$, it is expanded so as to spend $-\mathit{NC}$. Let $\MVPF_f\geq0$ be the MVPF of that policy at the margin, so that each dollar raised by cutting it costs its beneficiaries $\MVPF_f$ dollars of willingness to pay and each dollar spent on it delivers $\MVPF_f$ dollars. In both cases, the combined change raises welfare, in dollars, by
\begin{equation*}
\WTP^{\$}-\MVPF_f\,\mathit{NC}.
\end{equation*}
For $\mathit{NC}>0$, this is positive if and only if $\MVPF>\MVPF_f$. With $\mathit{FM}=\mathit{NC}/\mathit{ME}$ and $\mathit{ME}>0$, $\MVPF=(\WTP^{\$}/\mathit{ME})/\mathit{FM}$, so the condition is
\begin{equation*}
\frac{\WTP^{\$}}{\mathit{ME}}>\MVPF_f\times\mathit{FM}.
\end{equation*}
When the financing policy costs its payers one dollar of willingness to pay per dollar raised, $\MVPF_f=1$, and the condition compares the value of consumption smoothing with the fiscal multiplier, the two terms of the Baily--Chetty form used by \citet{haller2022welfare}. For $\WTP^{\$}>0$ and $\mathit{NC}\leq0$, the combined change raises welfare for every $\MVPF_f\geq0$; for $\mathit{ME}>0$, this is the case $\mathit{FM}\leq0$, in which the behavioral responses finance the entire mechanical cost.
